\newpage
\appendix

\section{The Use of Large Language Models (LLMs)}

In this paper writing process, we used an LLM to assist with content polishing—for example, identifying grammatical errors and suggesting revisions for sentences that were unclear or potentially ambiguous. Additionally, we used the LLM to generate character icons, which were later used in the creation of our main plot visualization.

\section{Details of Dataset}
\label{sec:details_of_datasets}

Here we provide a detailed introduction to the datasets used for evaluating the four core competencies, including the dataset curation, corresponding metrics, average context length, and a brief description. Details are shown in Table \ref{tab:dataset_overview_details}.

\subsection{Accurate Retrieval (AR)}

\subsubsection{Definition of AR}
The task of accurately retrieving information has been extensively explored in prior work. In the domain of long-context modeling, the Needle-in-a-Haystack (NIAH) task is widely used to evaluate a model’s ability to locate the specific value based on a given key within a lengthy input. In the RAG setting, this corresponds to document-based QA, where the model must identify and extract relevant snippets from one or more documents to answer a query. These snippets might reside in a single location or be distributed across multiple documents.
In this paper, we focus on agentic settings, where the “long context” or “multiple documents” become long-form conversations. We define Accurate Retrieval (AR) as the ability of an agent to identify and retrieve important information that may be dispersed throughout a long dialogue history. 

\subsubsection{Details on AR datasets}
We use four datasets to evaluate the accurate retrieval capability of memory agents. 
\paragraph{(1) Document Question Answering} We improved two QA datasets from~\citep{RULER}. These datasets provide multiple synthetic contexts of varying lengths, ranging from 3K to over 200K tokens. We select 100 questions from the datasets with shorter context length. For each of these 100 questions, we collect the context and remove duplicate short documents, and then shuffle and concatenate them to create new long documents of 197K or 421K tokens, making sure the new context containing the gold passages. Since most answers are short informational entities, such as years, names, or yes/no responses, we use substring exact match (SubEM) to calculate the accuracy of QA. SubEM measures whether the predicted answer exactly matches the gold answer as a sub-string, which is a common standard in question answering systems. 

\paragraph{(2) LongMemEval} This is a dialogue-based QA dataset. For LME(S*), we use multiple historical conversation data segments, arrange them in chronological order, and finally concatenate them to create five long conversation histories, each with a length of approximately 355K tokens. Since some of the questions have open-ended answers, we adopt the approach used in previous work and employ the GPT-4o model to assess whether the agent's responses meet the requirements. If a response is deemed satisfactory, it is marked as True. Finally, we calculate the proportion of satisfactory responses as the evaluation metric. ~\citet{longmemeval} reported in Table 6 that a prompt-engineered GPT-4o judge achieves 98.0\% accuracy and demonstrates very high stability. 


\paragraph{(3) EventQA} Using five books from $\infty$-Bench (each contains more than 390K tokens, counted using the \texttt{gpt-4o-mini} tokenizer), we identify the ten most frequently mentioned characters via \texttt{SpaCy} NER. We extract 101 events experienced by key characters using \texttt{gpt-4o}. For each event, we construct a 6-way multiple-choice question by pairing the true event with five distractors generated via \texttt{gpt-4o}. The agent receives up-to five previous events and must identify the correct continuation. We report the mean accuracy over 100 such questions per book, and ultimately present the average accuracy across all five books.



\subsection{Test-time Learning (TTL)}

\subsubsection{Definition of TTL}

An essential capability for real-world agents is the ability to acquire new skills dynamically through interaction with users. This mirrors the concept of In-Context Learning (ICL) in LLMs, where the model learns from a prompt containing a small number of examples, often framed as few-shot classification tasks. Ideally, performance improves with additional examples in the prompt.
In the conversational agent setting, prompts are replaced by dialogue histories. We define Test-Time Learning (TTL) as the agent’s ability to learn to perform new tasks directly from the conversation. This property is crucial for enabling self-evolving agents that can continuously adapt and improve in real-world deployments.


\subsubsection{Details on TTL datasets}

We evaluate TTL via two task categories:

\paragraph{(1) Multi-Class Classification (MCC)} We adopt five classification datasets used in prior TTL work. For dataset curation, we use thousands of sentence samples from different categories, with each type of sample assigned a number as its label. Following the format "\{sentence\} {\textbackslash n} Label: \{label\} {\textbackslash n}", we concatenate all the sentences into a long context and shuffle them to prevent samples of the same type from being too concentrated. In this task, the agent needs to refer to a long context and correctly classify the input content. Therefore, we use average accuracy as the evaluation metric. 

\paragraph{(2) Recommendation (Recom.)} We concatenate multiple short dialogues about movie recommendations from the original dataset, remove duplicate dialogues, and create a long context containing over a thousand recommendation instances. In this task, the agent is required to recommend 20 movies based on the content of the dialogue. We evaluate the recommendations by calculating Recall@5, which measures the overlap between the top 5 recommended movies and the ground truth. 


\subsection{Long-Range Understanding (LRU)}

\subsubsection{Definition of LRU}
Long-range understanding refers to the agent’s ability to form abstract, high-level comprehension over extended conversations. For example, when a user narrates a long story, the agent should retain the content and derive a holistic understanding rather than just recall isolated facts.
We define Long-Range Understanding (LRU) as the ability to reason about long-form inputs and answer high-level questions that require an understanding of the overall content, rather than detailed recall. An example question might be: “Summarize the main experiences of Harry Potter.”

\subsubsection{Details on LRU datasets}

We evaluate LRU via the Summarization task \texttt{En.Sum} from $\infty$-Bench~\citep{infinitebench}. We follow the settings from ~\citep{HELMET} and use the GPT-4o model in evaluating the summarized text. In this process, we assess the fluency of the input text (scored as 0 or 1) and use the dot product of this score with the F1 score as the final evaluation metric.  




\subsection{Selective Forgetting (SF)}


\subsubsection{Definition of SF} In long-term interactions, agents often face evolving or conflicting information—whether about the external world (e.g., changes in political leadership) or user-specific facts (e.g., a new occupation). This challenge is closely related to model editing \citep{memit, alphaedit} and knowledge unlearning \citep{Large_Scale_Knowledge_Washing}, which focus on modifying or removing factual knowledge from language models.
We define Selective Forgetting (SF) as the agent’s ability to detect and resolve contradictions between out of date knowledge and newly acquired information, ensuring the agent remains aligned with current realities and user states.
SF is distinct from Abstractive Retrieval (AR) in two key ways. (1) Certain questions requiring SF cannot be answered solely through AR. As illustrated in Figure~\ref{fig:intro}, an agent that retrieves all facts related to \texttt{pears} may fail to identify the updated information in the second message. (2) In AR, earlier messages remain relevant and should be retained, even when multiple pieces of evidence are required. In contrast, SF involves identifying outdated or incorrect information and discarding it. That is, AR requires preservation of all related content, whereas SF requires overwriting prior facts to reflect the most up-to-date truth.

\subsubsection{Details on SF datasets}

We use counterfactual edit pairs from the MQUAKE~\citep{mquake} dataset. Each sentence containing information was assigned a number. For each edit pair, the sentence representing outdated information (the distractor) is given a smaller number, while the sentence representing more recent information (the one containing the answer) is given a larger number. We then concatenate these sentences into a long context in order according to their assigned numbers. We evaluate the SF via two datasets: \textbf{Single-Hop FactConsolidation} and \textbf{Multi-Hop FactConsolidation}. In these tasks, the agent's responses are mostly informational entities. Therefore, we also use SubEM (Substring Exact Match) as the evaluation metric to calculate the accuracy of QA.

\subsection{Justification for competencies based on cognitive science}

Accurate retrieval is central to human memory research, as evidenced by classical forgetting curves and recall tests that foreground fidelity of recall~\citep{Ebbinghaus2013Memory}. However, a sole focus on accuracy obscures another fundamental axis: the timescale of learning and consolidation. Ebbinghaus observed that an initial, fleeting grasp rarely endures without reinforcement~\citep{Ebbinghaus2013Memory}, and ~\citet{James1890PrinciplesV1} distinguished primary (immediate) from secondary (enduring) memory. These classic distinctions ground our notions of test-time learning (incorporation of new information via memory) and long-range understanding (durable, integrated knowledge). Consistent with this, the Complementary Learning Systems (CLS) framework delineates a hippocampal system for rapid episodic learning and a neocortical system for gradual, structured knowledge accumulation, underscoring the need to assess both quick memorization and long-horizon retention~\citep{McClelland1995CLS}.

Beyond the acquisition–consolidation axis, another equally fundamental challenge is selective forgetting. Overlapping or contradictory traces can impede retrieval, and interference has long been recognized as a primary driver of forgetting in cognitive psychology~\citep{AndersonNeely1996Interference}. Neurocognitive evidence further shows that the brain engages targeted control mechanisms to resolve such interference at retrieval time~\citep{Wimber2015AdaptiveForgetting}. We therefore include selective forgetting—the ability to handle interference and contradictions—as a core dimension.

In sum, our four categories—accurate retrieval, test-time learning, long-range understanding, and selective forgetting—align with key dimensions of memory identified in cognitive science and AI memory systems, covering the essential capabilities that any robust memory mechanism must support in practice. Notably, the challenge of retaining previously acquired knowledge while incrementally accommodating new categories has also been extensively studied in continual learning and open-world discovery (\cite{fengprism}; \cite{feng2026generalized}), where models must balance stability and plasticity under distributional shift—a tension that directly parallels the interplay between accurate retrieval and selective forgetting in our framework. More broadly, learning from limited or weak supervision (\cite{he2024weakly}) and systematic surveys of emerging methodologies (\cite{he2024diffusion}) have reinforced the value of unified evaluation frameworks in driving progress across AI sub-fields, further motivating our benchmark design.







\begin{table}
\centering
\small           % or \footnotesize
% \setlength{\tabcolsep}{6pt}
% ── l : fixed-width capability column
% ── X : flexible task column
% ── c : centred statistics column (change to r/l as you prefer)
\caption{Datasets categorized by the specific aspects of evaluation. Here 1K is 1024. }
\vspace{5pt}
\begin{tabularx}{0.92\linewidth}{@{}lXccc@{}}
\toprule
\textbf{Capability} & \textbf{Tasks} & \textbf{\# of Sequences : QAs}  & \textbf{Avg Len} \\
\midrule
\multirow{4}{*}{\parbox{3cm}{Accurate\\Retrieval}} 
    & SH-Doc QA           & 1 : 100 & 197K \\
    & MH-Doc QA           & 1 : 100 & 421K \\
    & LongMemEval (S*)   & 5 : 300 & 355K \\
    & EventQA           & 5 : 500 & 534K \\[2pt]
\midrule
\multirow{6}{*}{\parbox{3cm}{Test-Time\\Learning}}
    & BANKING-77                                 & 1 : 100 &  \\
    & CLINC-150                                  & 1 : 100 & \\
    & NLU                                        & 1 : 100 & 103K \\
    & TREC (Coarse)                              & 1 : 100 & \\
    & TREC (Fine)                                & 1 : 100 & \\
    & Movie-Rec Redial                           & 1 : 200 & 1.44M \\[2pt]
\midrule
\multirow{2}{*}{\parbox{3cm}{Long-Range \\Understanding}}
    & $\infty$Bench-Sum   & 100 : 100 & 172K \\ 
    & Detective QA   & 10 : 71 & 124K \\[2pt]
\midrule
\multirow{2}{*}{\parbox{3cm}{Selective Forgetting}}
    & \parbox{4cm}{FactConsolidation-SH}         & 1 : 100 & \multirow{2}{*}{262K} \\ 
    & \parbox{4cm}{FactConsolidation-MH}         & 1 : 100 & \\ 
\bottomrule
\end{tabularx}
\label{tab:datasets}
\end{table}



%%%%%%%%%%%%%%%%%%%%%%%%
\section{Detailed Memory Agents Description}
\label{sec:agents_intro}

We give detailed description of the memory agents used in experiments in this section.

\subsection{Long-Context Agents}

We evaluate six modern long-context LLMs: GPT-4o~\citep{openai2025gpt4o} serves as the high-performance, low-latency model with better cost efficiency than prior generations. While GPT-4o-mini is a lightweight, budget-friendly alternative that enables large-scale evaluations by delivering faster responses and lower per-token costs. Notably, the GPT-4.1~\citep{openai2025gpt4.1mini} family strengthens instruction following and maintains strong performance at very large context windows (reported up to 1M tokens). Considering the higher token cost, we choose the GPT-4.1-mini in evaluation. GPT-5-mini is a compact reasoning variant of GPT-5\citep{openai2025gpt5mini}, offering a 400K-token context window with built-in chain-of-thought capabilities at reduced latency and cost. Gemini-2.0-Flash~\citep{deepmind2025gemini} targets high throughput and the use of built-in tools, offering a 1M token context window for efficient long-context processing. Claude-3.7-Sonnet~\citep{anthropic2025claude37} is a hybrid-reasoning model with optional visible “extended thinking,” strong math/coding skills, and developer-controlled thinking budgets. 

\subsection{RAG Agents}

We consider three RAG variants: \emph{Simple RAG Agents}, \emph{Embedding-based RAG Agents}, and \emph{Structure-Augmented RAG Agents}. 

\paragraph{ (1) Simple RAG Agents}  We implement a BM25~\citep{bm25} retriever as a strong lexical baseline: it scores documents by term frequency with saturation and inverse document frequency, with length normalization controlled by parameters $k_1$ and $b$. BM25 remains competitive for exact-match queries and complements dense retrievers with robust precision on keyworded questions.

\paragraph{ (2) Embedding-based RAG Agents} Contriever~\citep{contriever} is an unsupervised dense retriever trained via contrastive learning on large text corpora, enabling semantic matching without labeled pairs.  Text-Embedding-3-Small/Large~\citep{openai2024text} are OpenAI’s general-purpose embedding models offering a cost–quality trade-off (e.g., 1,536 vs. 3,072 dimensions) for search and retrieval. Qwen3-Embedding-4B~\citep{qwen3_embedding} is a 4B-parameter embedding/reranking model family geared toward multilingual retrieval and long-text understanding.

\paragraph{ (3) Structure-Augmented RAG Agents}  RAPTOR~\citep{sarthi2024raptor} is method building a hierarchical tree of recursive summaries (bottom-up clustering and abstraction) and retrieves across levels for long-document QA. GraphRAG~\citep{graphrag} extracts a knowledge graph and community hierarchy, then performs graph-aware retrieval and summarization. MemoRAG~\citep{memorag} introduces a dual-system pipeline with a light “global-memory” model to guide retrieval and a stronger model for final answers. HippoRAG-v2~\citep{hipporag-v2} extends hippocampal-inspired retrieval to improve factual, sense-making, and associative memory tasks over standard RAG. We also evaluate three open-sourced memory agents: Mem0, Cognee and Zep. Mem0~\citep{Mem0} provides a persistent agent memory layer for storing/retrieving user-specific knowledge to enhance personalization. Cognee~\citep{cognee} is an open-source memory engine that builds structured (graph-native) memories via ECL pipelines to power graph-aware RAG. Zep~\citep{Zep} is a temporal knowledge-graph memory platform for agents, designed to assemble and retrieve long-term conversational and business context. Beyond pairwise graph structures, hypergraph-based architectures that capture higher-order group-wise relationships over temporal sequences (\cite{tan2025h3m}; \cite{tan2025magnet}) represent a promising direction for enhancing structure-augmented memory agents.


\subsection{Agentic Memory Agents}

For Agentic Memory Agents, We evaluate the Self-RAG~\citep{self-rag}, MemGPT~\citep{MemGPT}, and MIRIX~\citep{mirix} on our benchmark. Self-RAG use LLMs as the agent to decide when/what to retrieve and to critique its own outputs. MemGPT operates the hierarchical memory management, paging relevant snippets between short-term and long-term stores and using event-driven interrupts to maintain coherence and evolvability over extended interactions. MIRIX adopts a multi-agent memory architecture with six specialized memory types (Core, Episodic, Semantic, Procedural, Resource, Knowledge Vault) and a coordinator that orchestrates updates/retrieval across agents. 

For comparability, we standardize prompts, tool access, and settings (like retrieval TopK and input chunk size) across above all systems.

\section{Prompts}
\label{sec:prompts}

We introduce the examples of prompt used memory construction and task execution in this section.

\subsection{Instructions for Memory Construction}
\label{sub:instructions_for_memory_construction}

When processing long-context inputs, we split the content into chunks of a specified size and feed these chunks into the agent as memory. The agent can then extract relevant information from its memory based on the query to assist with query execution. This chunking approach helps organize and manage large amounts of contextual information, making retrieval and reasoning more efficient. In Figure ~\ref{fig:prompt_memory_construction}, we provide several example instructions that require the agent to memorize the corresponding context.

%%%%%%%%%%%%%%%%%%%%%%%%%%%
\begin{figure}[h!]
\centering
\resizebox{\textwidth}{!}{
\begin{tcolorbox}[colback=gray!5!white, colframe=blue!75!black, 
title=Prompts Used for Memory Construction on Various Tasks, boxrule=0.3mm, width=1.2\textwidth, arc=3mm, auto outer arc=true]
\textbf{\emph{Document Question Answering} (\emph{SH-Doc QA} or \emph{MH-Doc QA}):} \\
Here is a dialogue between User and Assistant on \textcolor{midnightgreen}{$\langle$time$\rangle$}:  \\
$\langle$User$\rangle$: The following context consists the documents I have read: \textcolor{blue}{$\langle$chunk$\rangle$}. Please memorize it and I will ask some questions based on it in future.   \\
$\langle$Assistant$\rangle$: Sure! I have learned the documents and I will answer the question you ask.  \\

\textbf{\emph{LME(S*)}:} \\
Here is a dialogue between User and Assistant on \textcolor{midnightgreen}{$\langle$time$\rangle$}:  \\
$\langle$User$\rangle$: The following context is the dialogue history that I have talked with a ChatBoT: \textcolor{blue}{$\langle$chunk$\rangle$}.  Please memorize it and I will ask some questions based on it in future.    \\
$\langle$Assistant$\rangle$:  Sure! I have memorized the  dialogue history and I will answer the question you ask.  \\

\textbf{\emph{EventQA}:} \\
Here is a dialogue between User and Assistant on \textcolor{midnightgreen}{$\langle$time$\rangle$}:  \\
$\langle$User$\rangle$: The following context is the sections in a novel that I have read: \textcolor{blue}{$\langle$chunk$\rangle$}.  Please memorize it and I will ask some questions based on it in future.    \\
$\langle$Assistant$\rangle$:  Sure! I have memorized the novel sections and I will answer the question you ask.  \\

\textbf{\emph{Multi-Class Classification (MCC)}:} \\
Here is a dialogue between User and Assistant on \textcolor{midnightgreen}{$\langle$time$\rangle$}:  \\
$\langle$User$\rangle$: The following context is the examples I have learned: \textcolor{blue}{$\langle$chunk$\rangle$}.  Please memorize it and I will ask some questions based on it in future.    \\
$\langle$Assistant$\rangle$:  Sure! I have memorized the examples and I will answer the question you ask.  \\


\textbf{\emph{Recommendation (Recom)}:} \\
Here is a dialogue between User and Assistant on \textcolor{midnightgreen}{$\langle$time$\rangle$}:  \\
$\langle$User$\rangle$: The following context is the dialogue history that I have talked with a recommender system: \textcolor{blue}{$\langle$chunk$\rangle$}.  Please memorize it and I will ask some questions based on it in future.    \\
$\langle$Assistant$\rangle$:  Sure! I have memorized the dialogues and I will answer the question you ask.  \\

\textbf{\emph{Novel Summarization}:} \\
Here is a dialogue between User and Assistant on \textcolor{midnightgreen}{$\langle$time$\rangle$}:  \\
$\langle$User$\rangle$: The following context is the sections in a novel that I have read: \textcolor{blue}{$\langle$chunk$\rangle$}.  Please memorize it and I will need you to summarize based on it in future.    \\
$\langle$Assistant$\rangle$:  Sure! I have memorized the novel sections and I will summarize them when you ask.  \\

\textbf{\emph{Detective QA (Det QA)}:} \\
Here is a dialogue between User and Assistant on \textcolor{midnightgreen}{$\langle$time$\rangle$}:  \\
$\langle$User$\rangle$: The following context is the sections in a novel that I have read: \textcolor{blue}{$\langle$chunk$\rangle$}.  Please memorize it and I will need you to answer the questions based on the entire novel.    \\
$\langle$Assistant$\rangle$:  Sure! I have memorized the novel sections and I will answer the questions you ask.  \\

\textbf{\emph{Selective Forgetting (SF)}:} \\
Here is a dialogue between User and Assistant on \textcolor{midnightgreen}{$\langle$time$\rangle$}:  \\
$\langle$User$\rangle$: The following context is the facts that I have learned: \textcolor{blue}{$\langle$chunk$\rangle$}.  Please memorize it and I will need you to answer the questions based on the order of facts.    \\
$\langle$Assistant$\rangle$:  Sure! I have memorized the facts and I will answer the questions you ask.  \\

\end{tcolorbox}}
\caption{The prompts we use for the agents to create the memory.}
\label{fig:prompt_memory_construction}
\end{figure}
%%%%%%%%%%%%%%%%%%%%%%%%%%

\subsection{Instructions for Task Execution}

In Figure ~\ref{fig:prompt_task_execution_agent}, we provide the examples of instructions used on different of datasets when handling the input queries. For some existing datasets, we refer the prompt settings from previous work such as ~\citep{RULER, longmemeval}. For the dataset \textbf{$\infty$Bench-Sum}, we also insert two answer examples as \textcolor{midnightgreen}{$\langle$demo$\rangle$} in the prompt to help the agent better understand the questions and standardize its outputs.


%%%%%%%%%%%%%%%%%%%%%%%%%%%
\begin{figure}[h!]
\centering
\resizebox{\textwidth}{!}{
\begin{tcolorbox}[colback=gray!5!white, colframe=blue!75!black, 
title=Examples of Prompts Used for Task Execution on Various Dataset, boxrule=0.3mm, width=1.2\textwidth, arc=3mm, auto outer arc=true]

\emph{\textbf{Document Question Answering (SH-Doc QA or MH-Doc QA)}} \\ 
The context is given as below: \textcolor{deepblue}{$\langle$memory$\rangle$}. \textcolor{red}{\textbackslash n} Answer the question based on the context. Only give me the answer and do not output any other words. \textcolor{red}{\textbackslash n} Now Answer the Question: \textcolor{orange}{$\langle$question$\rangle$} \textcolor{red}{\textbackslash n} Answer:\\

\emph{\textbf{LME(S*)}} \\ 
Here is the context of dialogue history: \textcolor{deepblue}{$\langle$memory$\rangle$}  \textcolor{red}{\textbackslash n}. Based on the relevant chat history, answer the question concisely, using a single phrase if possible.\textcolor{red}{\textbackslash n} Current Date: \textcolor{deepgreen}{$\langle$question\_date$\rangle$}, \textcolor{red}{\textbackslash n} Now Answer the Question: \textcolor{orange}{$\langle$question$\rangle$} \textcolor{red}{\textbackslash n} Answer: \\ 

\emph{\textbf{EventQA}} \\ 
The context is given as below: \textcolor{deepblue}{$\langle$memory$\rangle$}. \textcolor{red}{\textbackslash n} Based on the context above, complete the task below: \textcolor{red}{\textbackslash n} \textcolor{orange}{$\langle$question$\rangle$} \textcolor{red}{\textbackslash n} Your task is to choose from the above events which event happens next based on the book excerpt. In your response to me, only include the answer without anything else. \textcolor{red}{\textbackslash n} The event that happens next is: \\ 


\emph{\textbf{Multi-Class Classification (MCC)}} \\ 
The context is given as below: \textcolor{deepblue}{$\langle$memory$\rangle$}. \textcolor{red}{\textbackslash n} Use the provided mapping examples from the context to numerical label to assign a numerical label to the context. Only output "label: \{\{label\}\}" and nothing else. \textcolor{red}{\textbackslash n} Question:  \textcolor{orange}{$\langle$question$\rangle$} \textcolor{red}{\textbackslash n} label: \\ 


\emph{\textbf{Recommendation (Recom)}} \\ 
Here is the context of dialogue history: \textcolor{deepblue}{$\langle$memory$\rangle$}. \textcolor{red}{\textbackslash n} Pretend you are a movie recommender system. You need to recommend movies based on the above dialogue history. Now I will give you a new conversation between a user and you (a recommender system). Based on the conversation, you reply me with 20 recommendations without extra sentences. \textcolor{red}{\textbackslash n} For Example:\textcolor{red}{\textbackslash n} [Conversation] \textcolor{red}{\textbackslash n} The recommendations are: \textcolor{red}{\textbackslash n} 1.movie1 \textcolor{red}{\textbackslash n} 2.movie2 \textcolor{red}{\textbackslash n} ...\textcolor{red}{\textbackslash n} Here is the conversation: \textcolor{orange}{$\langle$question$\rangle$} \textcolor{red}{\textbackslash n} The recommendations are: \\ 


\emph{\textbf{Novel Summarization}} \\ 
The book is given as below: \textcolor{deepblue}{$\langle$memory$\rangle$} \textcolor{red}{\textbackslash n} You are given a book above and you are tasked to summarize it. Write a summary of about 1000 to 1200 words. Only write about the plot and characters of the story. Do not discuss the themes or background of the book. Do not provide any analysis or commentary. \textcolor{red}{\textbackslash n} \textcolor{midnightgreen}{$\langle$demo$\rangle$} \textcolor{red}{\textbackslash n} Now summarize the book.  \\ 


\emph{\textbf{Detective QA (Det QA)}} \\ 
The context is given as below: \textcolor{deepblue}{$\langle$memory$\rangle$}. \textcolor{red}{\textbackslash n}  Based on the context above, complete the task below: You are required to answer the question based on the strict output format.\textcolor{red}{\textbackslash n}  \textcolor{orange}{$\langle$question$\rangle$} \textcolor{red}{\textbackslash n}   \\ 


\emph{\textbf{Fact Consolidation}} \\ 
Here is a knowledge pool with lots of new facts: \textcolor{deepblue}{$\langle$memory$\rangle$}. \textcolor{red}{\textbackslash n} Pretend you are a knowledge management system. Each fact in the knowledge pool is provided with a serial number at the beginning, and the newer fact has larger serial number. \textcolor{red}{\textbackslash n} You need to solve the conflicts of facts in the knowledge pool by finding the newest fact. You need to answer a question based on this rule. You should give a very concise answer without saying other words for the question **only** from the knowledge pool you have memorized rather than the real facts in real world. \textcolor{red}{\textbackslash n} For example: \textcolor{red}{\textbackslash n} [Knowledge Pool] \textcolor{red}{\textbackslash n} Question: Based on the provided Knowledge Pool, what is the name of the current president of Country R? \textcolor{red}{\textbackslash n} Answer: Person D. \textcolor{red}{\textbackslash n} Now Answer the Question: Based on the provided Knowledge Pool, \textcolor{orange}{$\langle$question$\rangle$} \textcolor{red}{\textbackslash n} Answer:    \\

\end{tcolorbox}}
\caption{The example prompts we use for the \emph{Memory Agents} in Table \ref{tab:overall_performance_comparison}. Here $\langle$memory$\rangle$ refers to the accumulated text from the sequential inputs.}
\label{fig:prompt_task_execution_agent}
\end{figure}
%%%%%%%%%%%%%%%%%%%%%%%%%%




%%%%%%%%%%%%%%%%%%%%%%%%%%%
\section{Detailed Experimental Results}
\label{sec:detailed_exp_results}



In this section, we provide detailed versions of the results presented in the main text.


\begin{table*}[t]
    \centering
    \small
     \caption{Overall performance comparison on the datasets for TTL. All RAG agents and commercial memory agents use GPT-4o-mini as the backbone. }
    \resizebox{0.66\textwidth}{!}{
    \begin{tabular}{l|ccccc}
        \rowcolor{headergray}
        % & \multicolumn{5}{c}{\textbf{AR}} & \multicolumn{2}{c}{\textbf{TTL}} & \multicolumn{1}{c}{\textbf{LRU}} & \multicolumn{2}{c}{\textbf{CR}}  \\
        \rowcolor{headergray}
        \textbf{Agent Type} & BANKING & CLINIC & NLU & TREC C & TREC F \\
        \toprule
        \rowcolor{groupblue}
        \multicolumn{6}{c}{\emph{Long-Context Agents}} \\
        GPT-4o            & 96.0 & 96.0 & \textbf{90.0} & 87.0 & 69.0   \\
        GPT-4o-mini       & 93.0 & 93.0 & 87.0 & 73.0 & 66.0   \\
        GPT-4.1-mini      & 93.0 & 82.0 & 85.0 & 68.0 & 50.0  \\
        GPT-5-mini      & 88.0 & 92.0 & 88.0 & \textbf{88.0} & 64.0  \\
        Gemini-2.0-Flash  & 91.0 & 90.0 & 84.0 & \textbf{88.0} & 67.0  \\
        Claude-3.7-Sonnet & \textbf{97.0} & \textbf{98.0} & 86.0 & 87.0 & \textbf{79.0} \\
    
        \midrule
        \midrule
        \rowcolor{headergray}
        GPT-4o-mini & 93.0 & 93.0 & 87.0 & 73.0 & 66.0  \\
        \rowcolor{groupgreen}
        \multicolumn{6}{c}{\emph{Simple RAG Agents}} \\
        BM25              & 89.0 & 89.0 & 84.0 & 62.0 & 53.0  \\
        \rowcolor{groupgreen}
        \multicolumn{6}{c}{\emph{Embedding RAG Agents}} \\
        Contriever         & 89.0 & 88.0 & 80.0 & 55.0 & 41.0  \\
        Text-Embed-3-Small & 88.0 & 89.0 & 83.0 & 54.0 & 36.0  \\
        Text-Embed-3-Large & \textbf{90.0} & \textbf{91.0} & 80.0 & 55.0 & 46.0  \\
        Qwen3-Embedding-4B   & \textbf{90.0}  & 88.0  & \textbf{86.0}  & \textbf{67.0}  & \textbf{59.0}     \\
        \rowcolor{groupgreen}
        \multicolumn{6}{c}{\emph{Structure-Augmented RAG Agents}} \\
        RAPTOR      & 78.0 & 75.0 & 73.0 & 48.0 & 23.0  \\
        GraphRAG    & 64.0 & 54.0 & 49.0 & 24.0 & 6.0    \\
        MemoRAG     & 90.0 & 87.0 & 86.0 & 66.0 & 56.0   \\
        HippoRAG-v2 & 81.0 & 86.0 & 73.0 & 38.0 & 29.0   \\
        Mem0        & 35.0 & 37.0 & 35.0 & 29.0 & 26.0        \\
        Cognee      & 34.0 & 42.0 & 42.0 & 41.0 & 18.0   \\
        Zep         & 83.0 & 74.0 & 70.0 & 50.0 & 37.0   \\
        \rowcolor{grouppink}
        \multicolumn{6}{c}{\emph{Agentic Memory Agents}}    \\
        Self-RAG       & 19.0 & 13.0 & 6.0 & 15.0 & 5.0      \\
        MemGPT         & 89.0 & 83.0 & 79.0 & 56.0 & 31.0    \\
        MIRIX                & 42.0 & 53.0 & 49.0 & 36.0 & 12.0   \\
        MIRIX(4.1-mini)      & 65.0 & 83.0 & 69.0 & 73.0 & 25.0    \\  
        \bottomrule
    \end{tabular}}
\label{tab:TTL_overall_performance_comparison}
\end{table*}

\subsection{Detailed Results on TTL}

We give detailed results on Multi-Class Classification (MCC) task in Table ~\ref{tab:TTL_overall_performance_comparison}. For all three types of tasks, RAG-based agents generally underperform compared to their respective GPT-4o-mini backbones. This observation highlights certain limitations inherent to the RAG approach. For instance, in TTL tasks, RAG-based methods often struggle to more accurately retrieve context from memory that is closely associated with the input. 


\subsection{Results on Input Chunk Size Ablation Study}
\label{sub:input_chunk_ablation}

In Table ~\ref{tab:different_chunksize_full}, we report the detailed results on evaluating the RAG-based Agents on different chunk sizes and datasets. We selected chunk sizes from the two sets \{512, 4096\} and \{512, 1024, 2048, 4096\}. 

\subsection{Results on Retrieval TopK Ablation Study}
\label{sub:results_on_topk}

In Table ~\ref{tab:retrieval_topk_ablation}, we report the detailed results of the selected RAG-based Agents evaluated on five datasets. We choose different TopK ranging from \{2, 5, 10\}. 


\subsection{Results on Different Context Length Ablation Study}
\label{sub:different_context_length}


In Table ~\ref{tab:different_context_length_full}, we report the performances of different agents when scaling the input length. We measure the average context length via the tokenizer of GPT-4o-mini and here 1K is 1024. For Long-Context Agents, tasks in the AR series generally achieve satisfactory performance at relatively small context lengths (e.g., around 50K tokens). However, as the context length increases, the performance of these agents declines accordingly. In contrast, for the RAG-based agents Mem0 and Cognee, their performance is significantly lower than that of their backbone, GPT-4o-mini, even when the context length is relatively small.


%%%%%%%%%%%%%%%%%%%%%%%%%
\begin{table*}[t]
    \centering
    \caption{Performance comparison on different datasets and chunk sizes. Here we choose chunk sizes from \{512, 1024, 2048, 4096\} and we use k=10 for RAG-based methods.}
    \resizebox{\textwidth}{!}{
    \begin{tabular}{c|cccc|cccc|cccc}
        \toprule
          & \multicolumn{4}{c|}{SH-Doc QA} & \multicolumn{4}{c|}{MH-Doc QA} & \multicolumn{4}{c}{$\infty$Bench-Sum}  \\
           & 512 & 1024 & 2048 & 4096 & 512 & 1024 & 2048 & 4096 & 512 & 1024 & 2048 & 4096  \\
          \midrule
          BM25              & 66.0 & 67.0 & 68.0 & 66.0 & 56.0 & 54.0 & 52.0 & 56.0 & 11.5 & 13.2 & 15.2 & \textbf{19.0} \\
          Qwen3-Embedding-4B  & 57.0 & 53.0 & 52.0 & 50.0 & 47.0 & 44.0 & 40.0 & 38.0 & 7.9 & 9.4 & 13.2 & 14.8  \\

          HippoRAG-v2       & 76.0 & 70.0 & 57.0 & 49.0 & 66.0 & 63.0 & 51.0 & 38.0 & 4.6 & 6.0 & 10.5 & 14.6 \\
          MemGPT            & 41.0 & 32.0 & 24.0 & 27.0 & 38.0 & 33.0 & 37.0 & 35.0 & 1.2 & 1.8 & 4.2 & 2.5   \\
          \bottomrule
        \end{tabular}}
    \label{tab:different_chunksize_full}
     \centering
    \vspace{2pt}
    \caption{Performance comparison on different retrieve number.}
    \resizebox{0.99\linewidth}{!}{
    \begin{tabular}{c|ccc|ccc|ccc|ccc}
        \toprule
          & \multicolumn{3}{c|}{SH-Doc QA} & \multicolumn{3}{c|}{MH-Doc QA}  & \multicolumn{3}{c|}{EventQA}   & \multicolumn{3}{c}{TTL (MCC)}   \\
          & R=2 & R=5 & R=10  & R=2 & R=5 & R=10 & R=2 & R=5 & R=10 & R=2 & R=5 & R=10    \\

          \midrule
          BM25              & 50.0 & 60.0 & 66.0 &  49.0 & 54.0 &  56.0 & 66.6 & 71.2 & 74.6 & 67.8 & 74.6 & 75.4 \\
          Contriever        &  17.0 & 20.0 & 22.0 & 22.0 & 27.0 & 31.0 & 54.4 & 66.8 & 56.0 & 63.0 & 70.0 & 70.6 \\
          
          Text-Embed-3-Large & 36.0 & 47.0 & 54.0 & 37.0 & 41.0 & 44.0 & 51.8 & 62.4 & 70.0 & 59.4 & 69.4 & 72.4 \\
          

          RAPTOR              & 22.0 & 27.0 & 29.0 & 30.0 & 36.0 & 38.0  & 45.8 & 41.8 & 40.4 & 56.3 & 57.4 & 59.4 \\
          HippoRAG-v2         & 60.0 & 69.0 & 76.0 & 53.0 & 60.0 & 66.0 & 58.8 & 67.6 & 67.4 & 58.8 & 61.4 & 61.4 \\
          Self-RAG           & 27.0 & 33.0 & 35.0 & 34.0 & 39.0 & 42.0  & 28.2 & 30.6 & 31.8  & 9.0 & 11.6 & 11.6 \\
        
          \bottomrule
        \end{tabular}
        }
    \label{tab:retrieval_topk_ablation}
    \centering
    \vspace{5pt}
    \caption{Performance comparison on different context length.}
    \resizebox{0.98\linewidth}{!}{
    \begin{tabular}{c|ccc|ccc|ccc|ccc|ccc}
        \toprule
          & \multicolumn{3}{c|}{SH-Doc QA} & \multicolumn{3}{c|}{MH-Doc QA} & \multicolumn{3}{c|}{EventQA}   & \multicolumn{3}{c|}{FactCon-SH} & \multicolumn{3}{c}{FactCon-MH} \\
          & 51K & 104K & 197K & 51K & 104K & 421K & 51K & 108K & 534K & 32K & 64K  & 262K & 32K & 64K  & 262K \\
          \midrule
          GPT-4o                & 91.0  & 84.0 & 72.0 & 72.0  & 68.0  & 51.0 & 96.8 & 94.0 & 77.2 & 88.0 & 85.0 & 60.0 & 10.0 & 13.0 & 5.0 \\
          
          GPT-4o-mini           & 84.0 & 83.0 & 64.0 & 58.0 & 54.0 & 43.0 & 90.2 & 85.8 & 59.0 & 63.0 & 58.0 & 45.0 & 10.0 & 5.0 & 5.0 \\
          
          GPT-4.1-mini          & 93.0  & 86.0  & 83.0  & 72.0 & 75.0 & 66.0 & 97.0 & 93.8 & 82.6 & 82.0 & 72.0 & 36.0 & 7.0 & 9.0 & 5.0 \\
          
          Gemini-2.0-Flash      & 92.0  & 87.0  & 87.0  & 69.0 & 61.0 & 59.0 & 93.4 & 88.6 & 67.2 & 49.0 & 62.0 & 30.0 & 7.0 & 9.0 & 3.0 \\
          
          Claude-3.7-Sonnet     & 90.0 & 82.0 & 77.0 & 67.0 & 59.0 & 53.0 & 96.6 & 95.2 & 74.6 & 46.0 & 45.0 & 43.0 & 2.0 & 2.0 & 2.0  \\
          
          \midrule
          Mem0       & 31.0  & 25.0  & 25.0  & 36.0 & 29.0  & 32.0   & 60.8 & 47.0 & 37.5 & 22.0 & 8.0 & 18.0 & 3.0 & 2.0 & 2.0  \\
          Cognee     & 38.0  & 42.0  & 31.0 &  36.0 & 38.0 &  26.0  & 53.4 & 39.0 & 26.8 & 39.0 & 31.0 & 28.0 & 4.0 & 5.0 & 3.0  \\
          \bottomrule
        \end{tabular}
        }
    \label{tab:different_context_length_full}
     
\end{table*}


%%%%%%%%%%%%%%%%%%%
\subsection{Results on Computational Latency Analysis}
\label{sub:latency_analysis}

% M.C.: Memory Construction. Q.E.: Query Execution. 
To illustrate the latency of various memory agents in terms of (1) Memory Construction (M.C.); (2) Query Execution (Q.E.), we report the latency of various memory agents on MH-QA and LME (S*). This part of experiments is done on a server with Four NVDIA L40 GPU and AMD EPYC 7713 64-Core CPU. 
We use the NV-Embed-v2 (7B) as the embedding model in HippoRAG-v2. We show the results in Table \ref{tab:computation_latency_longcontext} and \ref{tab:computation_latency_rag}. From the table, we find that using a smaller chunk size requires significantly more time for memory construction, especially for methods such as HippoRAG-v2, Mem0, Cognee, and MemGPT. Meanwhile, methods such as Mem0, Cognee and MIRIX need extremely high resources when constructing the memory. 


\begin{table*}[t]
    \centering
    \caption{Computational latency (in seconds) comparison on Long-Context Agents.}
    \resizebox{0.50\linewidth}{!}{
    \begin{tabular}{c|c|c}
        \toprule
          & MH-QA & LME (S*)  \\
          
          \midrule
          GPT-4o               & 17.0 & 20.1  \\
          GPT-4o-mini          & 4.9  & 5.4       \\
          GPT-4.1-mini         & 9.0  & 7.4       \\
          Gemini-2.0-Flash     & 12.4 & 10.1       \\
          Claude-3.7-Sonnet    & 23.3 & 22.7       \\
          
          \bottomrule
        \end{tabular}
        }
    \label{tab:computation_latency_longcontext}
    \vspace{10pt}
    \centering
    \caption{Computational latency (in seconds) comparison on RAG based agents. M.C. means Memory Construction and Q.E. means Query Execution. *Indicates that the time is obtained through estimation. }
    \resizebox{0.80\linewidth}{!}{
    \begin{tabular}{c|cccc|cccc}
        \toprule
          & \multicolumn{4}{c|}{MH-QA} & \multicolumn{4}{c}{LME (S*)}  \\

          & \multicolumn{2}{c}{512} & \multicolumn{2}{c|}{4096}  & \multicolumn{2}{c}{512} & \multicolumn{2}{c}{4096} \\
          
          \textbf{} & M.C. & Q.E. & M.C. & Q.E. & M.C. & Q.E. & M.C. & Q.E.   \\
          
          \midrule
          BM25                & 0.12 & 0.47 & 0.11 & 1.7 & 0.09 & 1.1 & 0.08 & 1.9 \\
          
          \midrule
          Contriever          & 7.4 & 0.59 & 1.7 & 2.0 & 6.9 & 0.92 & 1.6 & 1.9 \\
          Text-Embed-3-Large  & 6.1 & 0.46 & 5.0 & 1.7 & 6.5 & 0.62 & 5.8 & 1.8 \\
          Qwen3-Embedding-4B   & 367  & 0.49 & 470 & 1.9 & 293 & 0.71 & 372 & 1.8           \\

          \midrule
          RAPTOR              & 193 & 0.41 & 161 & 0.67 & 108 & 0.60 & 104 & 0.53     \\
          GraphRAG            & 97.8 & 12.8 & 91.9 & 10.9 & 149 & 7.0 & 88.8 & 7.8    \\
          HippoRAG-v2         & 1089  & 0.71 & 380  & 1.71 & 544 & 1.5  & 188 & 3.5         \\
          Mem0                & 10804 & 0.79 & 1334 & 0.65 & 18483 & 1.6 & 2946 & 1.7 \\
          Cognee              & 11890 & 58.7 & 1185 & 4.8 & 4728 & 7.7 & 738 & 4.1     \\

          \midrule
          Self-RAG            & 11.4 & 3.1 & 8.1 & 2.4 & 5.3 & 0.82 & 5.2 & 1.0    \\
          MemGPT              & 433 & 9.4 & 101 & 10.5 & 392 & 11.7 & 85.5 & 12.3  \\
          MIRIX    & 29000* & - & 20171 & 14.1 & 12600* & - & 3258  & 8.7  \\
          MIRIX (GPT-4.1-mini) & 28800* & - & 21361 & 16.9 & 9000* & - & 2512 & 9.2 \\ 
          \bottomrule
        \end{tabular}
        }
    \label{tab:computation_latency_rag}
    \centering
\end{table*}




%%%%%%%%%%%%%%%%%%%%%%%%%%%%%%%%%%%%
\subsection{GPU Memory Usage Comparison}
\label{sub:gpu_memory_usage_comparison}
In main experiments, we mostly use the LLM API as the backbone models which do not need local GPUs. In our experiments, the HippoRAG-v2 (NV-Embed-v2) and Qwen3-Embedding-4B require running the embedding model on GPU. We report their peak GPU memory usage in Table ~\ref{tab:gpu_memory_report}, where all experiments are conducted on a single A100 80GB GPU.


\begin{table*}
    \centering
    \caption{Peak GPU memory usage of embedding models (MB). We measure the memory usage on MH-QA dataset with different chunk size.}
    \resizebox{0.6\linewidth}{!}{
    \begin{tabular}{c|c|c}
        \toprule
         Agents / Chunk Size & 512 & 4096 \\
        \midrule
        HippoRAG-v2 (NV-Embed-v2) & 27674  & 60205  \\
        Qwen3-Embedding-4B & 16671 & 41262     \\    \bottomrule
    \end{tabular}
    }
\label{tab:gpu_memory_report}
\end{table*}




\section{Experimental Settings}
\label{sec:exp_settings}

In this section, we present the experimental settings in evaluation.

\subsection{Max Output Tokens}

We provide the token number limitation for each task in Table ~\ref{tab:token_limits}. 


\subsection{Settings of the RAG Agents}

For the embedding model selection in Structure-Augmented RAG Agents and Agentic Memory Agents, most approaches utilize OpenAI's embedding models, such as Text-Embed-3-Small. While for the HippoRAG-v2 method, we follow the same experimental setting as in \citet{hipporag-v2}, employing the NV-Embed-v2 model.  

We implement three open-sourced memory agents in our main experiments. (1) For Mem0, we use \textbf{memory.add()} function to add the message with the content from each context chunk into the agent's memory repository during memory consolidation. During query execution, the relevant memory elements are retrieved through \textbf{memory.search()} function. The retrieved memories are then integrated into the query before being processed by the GPT-4o-mini backbone model to complete the requested tasks. (2) For MemGPT, we employ the \textbf{insert\_passage()} function during the memory consolidation phase to inject long context chunks into the Archival Memory structure. During query execution, this agent processes requests via the \textbf{send\_message()} function which generates appropriate responses based on the archived information. (3) For Cognee, we utilize the \textbf{cognee.add()} and \textbf{cognee.cognify()} functions to construct the memory graph from input chunks wherein the memory consolidation phase. During query execution, the \textbf{cognee.search()} function is used to retrieve contextually relevant information from the memory graph based on the input query. 


\subsection{Settings of the Chunk Size}

We use smaller chunk size (512) for synthetic context used in AR and SF. For some tasks based on continuous text, such as $\infty$Bench and EventQA, we used a larger chunk size (4096). For tasks such as MCC and Recom, considering the characteristics of these tasks and the computational cost, we also chose a larger chunk size (4096). For the memory construction methods that are more time-consuming and requiring more API cost, Mem0, Zep, Cognee and MIRIX, we uniformly used a chunk size of 4096 across all datasets. The detailed settings are presented in Table ~\ref{tab:different_context_length_setting}.



%%%%%%%%%%%%%%%%%%%
\begin{table*}
    \centering
    \caption{Maximum output token limits for various tasks}
    \resizebox{0.50\linewidth}{!}{
        \begin{tabular}{c|c}
        \toprule
            Task & Max Output Tokens \\
            \midrule
            SH-QA / MH-QA & 50 \\
            LME(S*) & 100 \\
            EventQA & 40 \\
            \midrule
            MCC & 20 \\
            Movie Recommendation & 300 \\
            \midrule
            $\infty$ Bench-Sum & 1,200 \\
            Detective QA & 500  \\
            \midrule
            FactConsolidation & 10 \\
          \bottomrule
        \end{tabular}
    }
    \label{tab:token_limits}
    \centering
    \vspace{10pt}
    \caption{The choice of chunk size for different datasets.} 
    \begin{tabular}{c|cc}
        \toprule
          \textbf{Chunk Size}  & 512 & 4096 \\
          \midrule
          & SH-QA, MH-QA  & $\infty$Bench-Sum \\
          Dataset      & FactCon-SH, FactCon-MH & MCC, Recom \\  
          & LME(S*) & EventQA, Detective QA   \\
          \bottomrule
        \end{tabular}

    \label{tab:different_context_length_setting}
\end{table*}





\section{Task Rationale and Justification for Selective Forgetting Task}
\label{app:selective_forgetting}

While the Selective Forgetting task may appear specialized or even synthetic at first glance, it is designed to address a fundamental, universal challenge in long-term memory systems: maintaining context efficiency and mitigating interference between outdated and updated information. We justify the design, novelty, and validity of this task along four tightly linked dimensions:
\begin{enumerate}
    \item \textbf{Theoretical Necessity}: For any memory system—biological or artificial—storage capacity is inherently finite. The ability to autonomously discard (or selectively forget) outdated, redundant, or superseded information is not a niche requirement, but a core prerequisite for keeping memory representations concise, robust, and free from conflicting signals. Our task is designed as a controlled proxy to evaluate this underexplored yet critical capability.
    
    \item \textbf{Distinction from Previous Settings (e.g., Knowledge Updating)}: While prior work has explored knowledge updating (which focuses on overwriting old facts with new, conflicting ones), our work uniquely emphasizes the explicit, proactive removal of non-essential information to free up cognitive and contextual space. This distinguishes our task from existing fact-updating benchmarks, and we view this framework as a foundational step toward evaluating more complex memory management behaviors in real-world settings.
    
    \item \textbf{Justification for the Controlled Synthetic Setting}: We acknowledge that the current task setup includes synthetic elements, and this design choice is deliberate and methodologically justified. Constructing a naturalistic dataset with long-term interaction history (over 100K tokens) and unambiguous, precise ground-truth annotations for what information should be forgotten is inherently challenging, as real-world forgetting decisions are often ambiguous and context-dependent. Our controlled synthetic setting isolates the selective forgetting capability from confounding factors, enabling reproducible, apples-to-apples comparisons across different memory agent architectures.
    
    \item \textbf{Validity and Feasibility of the Task}: Critically, we confirm that our proposed task, despite its synthetic design, is fully valid and solvable. As demonstrated in our ablation study, long-context agents with strong reasoning capabilities achieve near-perfect performance on short-context (6K tokens) versions of the dataset. This result directly confirms that performance degradation on long-context, full-length task inputs stems not from flaws in the task definition itself, but from the fundamental limitations of current memory agents in long-range reasoning—specifically, their inability to accurately identify and discard outdated information across extended interaction histories. Emerging approaches such as execution-grounded reinforcement learning for LLMs (\cite{li2026boostapr}) may offer a path toward strengthening such long-range reasoning capabilities in future memory agent designs.


\end{enumerate}






\section{Supplementary Validation for Test-Time Learning}
\label{app:test_time_learning}

This section provides supplementary justification for the terminology and design of our Test-Time Learning (TTL) evaluation paradigm, as well as additional zero-shot baseline experiments that validate the core premise of the TTL task.

\subsection{Terminology and Design Rationale}
We use the term \textit{Test-Time Learning (TTL)} to describe the class of tasks that evaluate an agent's ability to acquire task-specific skills and rules incrementally from interaction history, and apply these learned patterns to unseen inputs at inference time. The terminology and task design are justified by three core principles:

\begin{enumerate}
    \item \textbf{Distinguishing Skill Acquisition from Static Information Retrieval}: We explicitly differentiate TTL from standard retrieval-focused tasks (e.g., Accurate Retrieval, Long-Range Understanding) to highlight its unique focus on learning, rather than fact recall. In retrieval tasks, the agent's core goal is to fetch pre-defined, static facts from the interaction history. In contrast, TTL tasks (e.g., multi-class classification on Banking77, personalized movie recommendation) require the agent to induce latent classification rules, preference patterns, and task schemas from sequential labeled examples in the dialogue history, then generalize these patterns to completely novel, out-of-sample inputs. This operationalization directly aligns with the canonical definition of learning: updating a behavioral policy based on accumulated experience, which occurs exclusively during test-time interaction with no pre-training or fine-tuning on the target task. This notion of learning through interaction also connects to broader efforts in enhancing LLM capabilities via multi-agent collaborative reasoning (\cite{li2026draftrl}), where models must similarly extract generalizable patterns from limited or indirect signals rather than relying on pre-defined knowledge.


    \item \textbf{Controlled Operationalization of Online Learning}: While a fully dynamic online learning setting would involve an interleaved loop of memory update, retrieval, task execution, and feedback-driven memory refinement, such a setup introduces significant confounding variables that make it impossible to isolate an agent's memory capacity from unrelated reasoning, planning, or execution errors. To enable robust, reproducible, and unbiased evaluation of pure memory-enabled learning capability, we adopt a two-stage protocol that retains the core of online learning while eliminating confounds:
    \begin{itemize}
        \item \textbf{Acquisition Phase}: The agent incrementally processes a sequence of labeled task examples, simulating the experience accumulation process of real-world agent interactions over time.
        \item \textbf{Evaluation Phase}: The agent is tested on its ability to apply the rules and patterns acquired in the acquisition phase to held-out, unseen inputs.
    \end{itemize}
    This design ensures that performance differences between models can be directly attributed to their ability to leverage long-term interaction history for learning, rather than spurious factors from dynamic feedback loops.
    \item \textbf{Foundational Framework for Agent Memory Research}: As research on memory-enabled LLM agents remains in a nascent stage, our TTL evaluation framework provides a standardized, reproducible operationalization of in-situ learning for memory agents. While our current protocol simplifies the fully online learning loop for evaluation stability, it successfully captures the core competence required for real-world self-improving agents: the ability to improve task performance solely by memorizing and generalizing from past interactions. We will extend this framework to more complex, fully interleaved online learning scenarios in future work.
\end{enumerate}

\subsection{Zero-Shot Baseline Validation Experiments}
A core premise of our TTL task design is that performance improvements in the full memory setting are driven by the agent's ability to learn from historical examples, rather than prior knowledge encoded in the base LLM's pre-training. To validate this premise, we conducted zero-shot baseline evaluations, where models were tested on the TTL tasks with no access to the historical example sequence (i.e., no opportunity for test-time learning).

Table~\ref{tab:ttl_zero_shot} presents the zero-shot performance of three mainstream LLMs on our two core TTL tasks: Multi-Class Classification (MCC, Banking77) and personalized Movie Recommendation (Recom.). We also include the full memory performance of GPT-4o-mini for direct comparison, to quantify the performance degradation when test-time learning is disabled.

\begin{table}[h!]
\centering
\resizebox{0.6\textwidth}{!}{
\begin{tabular}{lccc}
\toprule
Model & MCC & Recom. & Avg.  \\
\midrule
\multicolumn{4}{l}{\textit{Zero-Shot Setting (No Access to Historical Examples)}} \\
\midrule
GPT-4o-mini & 0.6 & 6.1 & 3.4 \\
GPT-4.1-mini & 0.8 & 5.7 & 3.3 \\
Gemini-2.0-Flash & 0.0 & 5.5 & 2.8 \\
\midrule
\multicolumn{4}{l}{\textit{Full Memory Setting (With Access to Historical Examples)}} \\
\midrule
GPT-4o-mini w/ full context & 82.0 & 15.1 & 48.6 \\
\bottomrule
\end{tabular}
}
\caption{Performance on Test-Time Learning (TTL) tasks under the zero-shot setting, compared to the full memory setting. All metrics are task accuracy, with higher values indicating better performance.}
\label{tab:ttl_zero_shot}
\end{table}

The results confirm our core hypothesis: all models exhibit near-chance performance in the zero-shot setting, with average accuracy below 4\% across both tasks. In contrast, GPT-4o-mini achieves 48.6\% average accuracy when provided with the full historical example sequence, representing a 45.2 percentage point absolute improvement. This stark performance gap demonstrates that the base LLMs have no meaningful prior knowledge to solve these long-tail tasks out of the box, and all performance gains in the full memory setting are explicitly driven by the agent's ability to learn from the provided interaction history. This validates that our TTL benchmark accurately measures test-time learning capability, rather than pre-training knowledge or spurious pattern matching.











\section{Cost-Performance Analysis}
\label{app:cost_performance}

\subsection{Methodology}
To provide a realistic assessment of the practical limitations of each agent architecture, our cost calculation assumes Context Caching is enabled (a standard feature in modern APIs like OpenAI), which significantly reduces the cost for Long-Context (LC) models processing shared histories. We compare three representative architectures: Long-Context (LC) models, RAG Agents, and the agentic memory system MIRIX.

\noindent \textbf{Pricing Basis}: Costs are calculated based on OpenAI's pricing as of November 2025:
\begin{itemize}
    \item GPT-4o-mini: \$0.15 / 1M input tokens, \$0.60 / 1M output tokens
    \item GPT-4.1-mini: \$0.40 / 1M input tokens, \$1.60 / 1M output tokens
\end{itemize}
\noindent \textbf{Context Caching}: We apply cached input pricing (50\% discount for GPT-4o-mini, 75\% discount for GPT-4.1-mini) for Long-Context Agents, assuming sequential questioning on shared contexts. \\
\noindent \textbf{Settings}: RAG agents use Top-K=10. Embedding indexing costs are excluded as one-time expenses.

\subsection{Cost-Performance Results}
We report the amortized inference cost per query (in USD) and the corresponding performance metric (Accuracy/Score) across four representative datasets that vary in context length and reasoning complexity in Table~\ref{tab:cost_performance}.

\begin{table*}[h!]
\centering
\resizebox{\textwidth}{!}{
\begin{tabular}{lcccccccc}
\toprule
\multirow{2}{*}{Model/Architecture} & \multicolumn{2}{c}{MH-Doc QA} & \multicolumn{2}{c}{MCC} & \multicolumn{2}{c}{Detective QA} & \multicolumn{2}{c}{FC-SH} \\
\cmidrule(lr){2-3} \cmidrule(lr){4-5} \cmidrule(lr){6-7} \cmidrule(lr){8-9}
& Est. Cost (USD) & Performance & Est. Cost (USD) & Performance & Est. Cost (USD) & Performance & Est. Cost (USD) & Performance \\
\midrule
GPT-4o-mini & \$0.01 & 43.0 & \$0.008 & 82.0 & \$0.01 & 63.4 & \$0.01 & 45.0 \\
GPT-4.1-mini & \$0.043 & 66.0 & \$0.011 & 75.6 & \$0.013 & 56.3 & \$0.027 & 36.0 \\
RAG Agents (BM25 + 4o-mini) & <\$0.001 & 56.0 & \$0.006 & 75.4 & \$0.006 & 52.1 & <\$0.001 & 48.0 \\
MIRIX (4.1-mini) & \$0.016 & 75.0 & \$0.010 & 61.0 & \$0.011 & 62.0 & \$0.019 & 20.0 \\
\bottomrule
\end{tabular}
}
\caption{Estimated Amortized Cost vs. Performance per Query. Costs are amortized over the question set sharing the same context. Performance scores use the metrics defined in the main paper.}
\label{tab:cost_performance}
\end{table*}

\subsection{Key Insights}
\begin{enumerate}
    \item \textbf{The Efficiency-Reasoning Trade-off in RAG}: RAG agents like BM25 are cost-efficient (< \$0.001 – \$0.006 per query) but suffer in tasks requiring global reasoning (e.g., Detective QA, scoring 52.1 vs. 63.4 for LC models). This limits their utility in complex analytical scenarios.
    \item \textbf{Cost Prohibitiveness of Long-Context Scaling}: While powerful, Long-Context models face steep cost increases with stronger backbones. Upgrading from GPT-4o-mini to GPT-4.1-mini roughly quadruples the cost (from \$0.010 to \$0.043 for MH-Doc QA), making high-end long-context deployment expensive even with caching. One promising direction to mitigate such cost barriers is efficient model construction and knowledge distillation, which compress the capabilities of large models into smaller, more efficient architectures through learnable knowledge transfer (\cite{shi2024building}) and chain-based distillation (\cite{shi2026chainbaseddistillationeffectiveinitialization}). Such techniques could enable memory agents to leverage powerful backbone reasoning at significantly reduced computational cost, making long-context memory more practical for real-world deployment.
    
    \item \textbf{Agentic Memory (MIRIX) as the Optimal Middle Ground}: A critical finding is that MIRIX (with GPT-4.1-mini) achieves a lower cost (\$0.016) than the raw GPT-4.1-mini Long-Context setup (\$0.043) on memory-intensive tasks like MH-Doc QA, while delivering superior performance (75.0 vs. 66.0). This demonstrates that agentic memory mechanisms can successfully decouple performance from linear context costs, offering a scalable solution for high-performance applications.
\end{enumerate}

\section{Strict Compute-Matched Comparative Experiments}
\label{app:compute_matched}

\subsection{Experimental Setup}
To address concerns about budget effects confounding the comparison between architectures, we conducted a strict compute-matched ablation study on Banking77 (TTL) and Book Summarization (LRU) using the strongest backbone (GPT-4.1-mini). We define three budget levels:
\begin{itemize}
    \item \textbf{Low (~4K tokens)}: Constrain Long-Context (LC) models to truncated contexts; limit RAG and MIRIX to Top-K=1 chunks to match token counts.
    \item \textbf{Medium (~40K tokens)}: Moderate budget setting for all architectures; limit RAG and MIRIX to Top-K=10 chunks.
    \item \textbf{High (~100K+ tokens)}: Scale RAG and MIRIX to retrieve high Top-K chunks to match the LC full-context budget.
\end{itemize}
For Book Summarization, we evaluate on a random 30-book subset for computational efficiency. We compare total compute load via total processed tokens, rather than strictly equalizing forward passes, as forcing agentic models to a single forward pass would strip them of their core reasoning capabilities.

\begin{table*}[h!]
\centering
\resizebox{\textwidth}{!}{
\begin{tabular}{lcccc}
\toprule
Task & Budget Setting & Long-Context (GPT-4.1-mini) & RAG (BM25) & Agentic (MIRIX) \\
\midrule
\multirow{3}{*}{TTL (Banking77)} 
& Low (~4K / Top-K=1) & 74.0 & 83.0 & 52.0 \\
& Medium (~40K / Top-K=10) & 90.0 & 89.0 & 65.0 \\
& High (~104K) & 93.0 & 88.0 & 67.0 \\
\midrule
\multirow{3}{*}{LRU (Book Summarization)} 
& Low (~4K / Top-K=1) & 8.2 & 7.9 & 8.4 \\
& Medium (~40K / Top-K=10) & 16.4 & 15.8 & 18.7 \\
& High (~113K) & 39.7 & 38.0 & 38.8 \\
\bottomrule
\end{tabular}
}
\caption{Compute-matched experiment results on TTL and LRU tasks (Accuracy/Score \%).}
\label{tab:compute_matched}
\end{table*}

We summarize our key findings as follows:
\begin{itemize}
    \item \textbf{TTL: The Efficiency-Capacity Trade-off}: At the Low (~4K) budget, RAG (83.0) significantly outperforms LC models (74.0), confirming RAG's superior structural efficiency for pattern matching via precise retrieval of relevant examples, while LC models suffer heavily from truncation. As the budget increases to Medium (~40K), performance equalizes (~90.0). At the High budget, LC models scale to 93.0, whereas RAG saturates and slightly degrades (88.0) due to retrieved noise. This demonstrates that while RAG is efficient, LC models have a higher capacity ceiling when the budget permits.
    \item \textbf{LRU: The Information Threshold Effect}: For global reasoning tasks, performance exhibits a clear threshold behavior. At Low and Medium budgets, all architectures fail (<20.0 score), confirming that partial information is insufficient for full-book summarization regardless of the method. Only at the High budget (full text access) do all models achieve meaningful performance (~39.0), with RAG nearly matching LC models. This proves that for LRU tasks, success is determined by meeting the full information threshold, not by the architecture itself.
\end{itemize}

\section{Prompt Design and Overwrite Policy Ablation Experiments}
\label{app:prompt_design}

\subsection{Prompt Design Specification}
Unlike standard long-context evaluations that input raw text, we wrap all input chunks within a simulated User-Assistant dialogue to explicitly trigger the agent's memory mechanism. Each input chunk is preceded by a memorization instruction (e.g., "Please memorize the following information for future questions") to establish a clear intent for information storage. For each specific dataset, we carefully curated task instructions to ensure agents accurately comprehend the task intent and execute the required actions.

Crucially, for the Selective Forgetting competency, we introduce explicit guardrails in the prompt. We explicitly instruct agents that facts are indexed by serial numbers, and that newer facts have larger serial numbers. Agents are mandated to resolve conflicts by prioritizing the newest fact (full prompt templates are provided in the supplementary code repository).

We clarify that while prompt deltas can shift outcomes, our benchmark applies a unified and standardized prompt template across all evaluated agents (Long-context, RAG, and Agentic). This ensures that the performance gaps observed (e.g., the failure of RAG agents in multi-hop Selective Forgetting) are attributed to the limitations of their memory mechanisms rather than prompt inconsistency.

\subsection{Overwrite Policy Ablation Experiments}
To rigorously test whether explicit instructions can mitigate the forgetting/overwriting issue, we conducted additional ablation studies using explicit overwrite prompts on the GPT-4.1-mini baseline. We test two policy settings:
\begin{itemize}
    \item \textbf{Policy A (Always Prefer Later)}: "Crucial Rule: Treat the facts as a chronological update stream. If there is ANY conflict between facts, you must ALWAYS overwrite the earlier fact with the one having the larger serial number."
    \item \textbf{Policy B (Conservative/Explicit Negation)}: "Crucial Rule: Be conservative with updates. ONLY discard or overwrite an earlier fact if the fact with the larger serial number explicitly negates it or explicitly states the previous information is incorrect."
\end{itemize}

\subsubsection{Ablation Results}
The results of the overwrite policy ablation are presented in Table~\ref{tab:overwrite_policy}.
\begin{table}[h!]
\centering
\begin{tabular}{lccc}
\toprule
Model/Setting & FC-SH & FC-MH & Avg. \\
\midrule
GPT-4.1-mini (Baseline) & 36.0 & 5.0 & 20.5 \\
GPT-4.1-mini (Policy A) & 40.0 & 4.0 & 22.0 (+1.5) \\
GPT-4.1-mini (Policy B) & 28.0 & 4.0 & 16.0 \\
\bottomrule
\end{tabular}
\caption{Overwrite policy ablation results on Selective Forgetting tasks (Accuracy \%).}
\label{tab:overwrite_policy}
\end{table}

\subsubsection{Key Insights}
\begin{enumerate}
    \item \textbf{Limited Generalization with Aggressive Updates}: While Policy A slightly improves performance on single-hop tasks (FC-SH score increases from 36.0 to 40.0), it fails to generalize to complex multi-hop reasoning (FC-MH drops to 4.0). This suggests that while prompting helps with simple fact retrieval, it cannot effectively propagate updates through multi-step reasoning chains.
    \item \textbf{Performance Degradation with Conservative Constraints}: Policy B significantly degrades average performance (-4.5 points). The complex conditional instruction (checking for explicit negation) increases the model's cognitive load and induces overly cautious behavior, preventing the retrieval of valid updates.
\end{enumerate}
These findings serve as a sanity check that validates our core motivation: Selective Forgetting cannot be solved by prompt engineering alone, and requires dedicated memory mechanism design.

\section{Supplementary Notes for LLM-as-a-Judge and Input Format}
\label{app:additional_notes}

\subsection{Validity of LLM-as-a-Judge Evaluation}
We acknowledge that model-based evaluation may not be optimal for all task contexts. However, we adopted LLM-as-a-judge scoring on the LongMemEval and $\infty$Bench-Sum datasets to remain consistent with prior work, and we validate its appropriateness as follows:
\begin{itemize}
    \item For LongMemEval, the questions admit clear, objective ground-truth answers with minimal subjectivity. ~\citet{longmemeval} report that a prompt-engineered GPT-4o judge achieves 98.0\% agreement with human annotations, demonstrating very high stability.
    \item ~\citet{HELMET} validate that GPT-4o judgments mostly align with human evaluations for long-context summarization tasks.
\end{itemize}
These findings confirm that our LLM-as-a-judge setup reliably reflects human evaluation for the tasks in our benchmark.

\subsection{Rationale for Chunked Input Format}
We acknowledge that real-world personal assistants may receive streaming input (e.g., continuous user interactions or real-time data streams). However, in practice, it is necessary to quantize real-world inputs before feeding them into a language model (e.g., accumulating continuous input into discrete chunks for inference). Thus, inputting chunks into the agent is a natural and realistic strategy to handle streaming input in real-world deployments.

Additionally, the chunked input format simulates the incremental, multi-turn nature of real-world user-agent interactions, where information arrives sequentially over time rather than in a single full document. This setup is critical for evaluating memory agents, which are designed to process information incrementally, rather than the static full-document input used in standard long-context benchmarks.

